\chapter{Architektur}
\label{chap:architektur}

Die vorliegende Architektur folgt dem Paradigma der \gls{gls-NFV}. Das \gls{gls-ETSI} definiert dieses Paradigma als Standard. Das Ziel davon ist es, spezialisierte Hardware durch software- und containerbasierte Dienste zu ersetzen. Diese Dienste sollen stattdessen auf Standard-Servern betrieben werden können \parencite{ETSI-NFV}. Auch softwarebasierte \gls{gls-PBX}- und Telefonsysteme lassen sich im Sinne des \gls{gls-NFV}-Modells als \gls{gls-VNF} einordnen, da sie zentrale Kommunikationsfunktionen wie Signalisierung, Vermittlung und Medienverarbeitung übernehmen und virtualisiert betrieben werden können. Damit wird die Flexibilität und Skalierbarkeit klassischer Kommunikationssysteme anhand der Intention des \gls{gls-ETSI}-Referenzmodells erhöht.

\section{\gls{gls-ETSI} \glstext{gls-NFV}-Architektur-Standard}
\label{sec:etsi_nfv_architektur}

Die \gls{gls-ETSI}-\gls{gls-NFV}-Architektur bietet ein flexibles Rahmenwerk für die Virtualisierung von Netzwerkfunktionen auf Standardhardware und bildet damit die Grundlage für skalierbare Kommunikationssysteme. Sie unterscheidet drei zentrale Schichten: Die Infrastruktur stellt virtualisierte Ressourcen wie Rechenleistung, Speicher und Netzwerk bereit. Darauf aufbauend werden einzelne Netzwerkfunktionen als Softwaremodule betrieben, die unabhängig voneinander auf dieser Infrastruktur laufen. Die Steuerung und Verwaltung aller Komponenten übernimmt die \gls{gls-MANO}-Eben, die aus spezialisierten Modulen für das Lebenszyklusmanagement der \gls{gls-VNF}, der Ressourcenverwaltung und der übergreifenden Koordination besteht \parencite{ETSI-NFV}. Diese modulare Architektur ermöglicht es, spezialisierte Dienste wie Signalisierung, Medienverarbeitung und \gls{gls-KI}-Funktionen als unabhängig skalierbare Services zu implementieren. Unterstützt wird dies durch Orchestrierung und Failover-Mechanismen \parencite{ETSI-REL}.

Der Prototyp dieser Arbeit implementiert ein minimales Subset dieser Architektur mit Single-Host-Deployment statt verteilter Infrastruktur. Alle produktiven \gls{gls-ETSI}-Standards werden als optionale Upgrade-Strukturen angesehen und sind in hochskalierten Nutzungsszenarien erforderlich.

\subsection{Architektur-Übersicht}

Das konkrete System in dieser Arbeit realisiert einen durchgängigen Kommunikationspfad zwischen klassischer Telefonie und \gls{gls-KI}-basierter Sprachverarbeitung. Ein eingehender Anruf wird vom \gls{gls-SIP}-Gateway entgegengenommen, die Medienschicht transkodiert das Audio zwischen verschiedenen Codecs wie G.711, Opus und PCM16 und die Anwendungsschicht verarbeitet Sprache durch \gls{gls-ASR}, \gls{gls-LLM} und \gls{gls-TTS}. Die generierte Sprachantwort durchläuft den gleichen Pfad in umgekehrter Richtung zurück zum Anrufer, wobei jede Schicht isoliert, skaliert und gesichert werden kann.

Ein aktuelles Beispiel für eine moderne, \gls{gls-KI}-basierte Telefoniearchitektur liefert das System von Todorov, Todor N, das für die Automatisierung von Taxi-Bestellungen per Telefon entwickelt wurde. Die Architektur besteht aus einem offenen \gls{gls-PABX} für das Call-Routing, Echtzeit-\gls{gls-STT}/\gls{gls-TTS}-Modulen, einem domänenspezifisch trainierten \gls{gls-LLM}-Kern und robusten Schnittstellen zu Legacy-Dispatch-Systemen. Diese Struktur ist darauf ausgelegt, hohe Anrufvolumina mit minimaler Latenz zu verarbeiten und bietet Fallback-Mechanismen zu menschlichen Operatoren \parencite{Todorov2025LLMDispatcher}.

\subsection{Einordnung anhand von Architektur-Schichten}

Die in dieser Arbeit verwendete Architektur gliedert sich in drei logische Schichten und wird wie folgt beschrieben:

Die \textbf{Signalisierungsschicht}, die für die \gls{gls-SIP}-Steuerung zuständig ist, die \textbf{Medienschicht}, die \gls{gls-RTP}-Audio verwaltet, und die \textbf{Anwendungsschicht} übernehmen die Steuerung der \gls{gls-KI}-Logik. Diese Schichtung ermöglicht erweiterbare und isolierte Komponenten nach dem \gls{gls-NFV}-Paradigma \parencite{ETSI-NFV}.

\textbf{Signalisierungsschicht:}\\
Das \gls{gls-SIP}-Gateway \parencite{rfc3261} empfängt INVITE-Requests über \gls{gls-UDP})sowie \gls{gls-TCP} Port 5060 und führt \gls{gls-NAT}-Traversierung mit erweiterten Funktionen wie \gls{gls-STUN}/\gls{gls-TURN}/\gls{gls-ICE} durch \parencite{rfc5389, rfc5766, rfc5245}. Ein Applikationsserver verarbeitet \gls{gls-SIP}-Signale und erstellt Room-Sessions als Mikroservice.
\clearpage
\textbf{Medienschicht:}\\
Ein \gls{gls-RTP}-Relay leitet \gls{gls-RTP}-Pakete vom \gls{gls-SIP}-Telefon zum \gls{gls-WebRTC}-Server weiter und transkodiert die notwendigen Audio Codecs (G.711 $\leftrightarrow$ Opus $\leftrightarrow$ PCM16), die in RFC 3550 standardisiert sind \parencite{rfc3550}. Der \gls{gls-WebRTC}-Server terminiert \gls{gls-WebRTC}-Streams und stellt Opus-Audio mit 48 kHz bereit. Produktive Systeme implementieren zusätzlich \gls{gls-SRTP}-Verschlüsselung nach RFC 3711 \parencite{rfc3711}.

\textbf{Anwendungsschicht:}\\
Die Anwendungsschicht integriert \gls{gls-KI}-Dienste für Sprachverarbeitung und Dialog-Management. Sie verbindet Audio-Streams mit \gls{gls-ASR} für Spracherkennung, \gls{gls-LLM} für Intent-Verarbeitung und Antwort-Generierung sowie \gls{gls-TTS} für Sprachsynthese. \gls{gls-VAD} ermöglicht dabei die Erkennung von Sprachaktivität zur Optimierung der Verarbeitungskette. Diese Schicht realisiert den interaktiven \gls{gls-KI}-Dialog zwischen Anrufer und System.\parencite{Antonova2025,rfc3551,rfc3550}

\begin{table}[H]
\centering
\caption[Zielarchitektur-Komponenten]{Zielarchitektur: Komponenten und Funktionen nach Antonova et al., \gls{gls-ETSI}-\gls{gls-NFV}-Richtlinien und Todorov \parencite{Antonova2025, ETSI-REL, Todorov2025LLMDispatcher}}
\label{tab:vnf_components}
\small
\begin{tabular}{|>{\raggedright\arraybackslash}p{3.6cm}|>{\raggedright\arraybackslash}p{5.3cm}|>{\raggedright\arraybackslash}p{4.1cm}|}
\hline
\textbf{Komponente} & \textbf{Funktion} & \textbf{Charakteristik} \\
\hline
\gls{gls-PABX} (Asterisk) & Eingehende Anrufe annehmen, Audio an \gls{gls-STT} leiten, Call-Rerouting & Hochvolumige Anrufverwaltung, geringe Latenz \\
\hline
\gls{gls-SIP}-Proxy (Kamailio) & \gls{gls-SIP}-Nachrichten-Routing, Lastverteilung vor der \gls{gls-PABX} & Entlastung des \gls{gls-PABX}-Kerns \\
\hline
\gls{gls-RTP}-Proxy & \gls{gls-RTP}-Relay, Medienpfad unterstützen & Entkopplung des Medienpfads \\
\hline
\gls{gls-STT}/\gls{gls-TTS} & Sprache $\rightarrow$ Text, Text $\rightarrow$ Sprache & Echtzeit-Verarbeitung, robuste Spracherkennung \\
\hline
\gls{gls-LLM}-Kern mit Task-Orchestrierung & Dialogsteuerung, Kontextverarbeitung, \gls{gls-API}-Aufrufe & Niedrige Inferenzlatenz, domänenspezifisches Training \\
\hline
Externe \gls{gls-API}-Schnittstellen & Geocodierung, Adressvalidierung, Dispatch-Integration & Interoperabilität mit Legacy-Systemen \\
\hline
Fallback-Mechanismen & Rerouting zu menschlichen Operatoren & Servicekontinuität bei Fehlerfällen \\
\hline
Datenbank (PostgreSQL/MySQL) & Konfiguration, \gls{gls-CDR}, Authentifizierung & Zuverlässige Datenhaltung \\
\hline
Konfigurationsmanagement (ConfigMaps/Geheimnisse) & Zentrale Verwaltung von Konfigurationen und Geheimnisse & Trennung von Code und Geheimnissen \\
\hline
Load Balancer/Ingress & Traffic-Verteilung, Cluster-Einstieg & Skalierbarer Zugriffspunkt \\
\hline
Monitoring/Logging (Prometheus/Grafana) & Metriken und Log-Analyse & Betriebsüberwachung und Fehleranalyse \\
\hline
Skalierbare Architektur \gls{gls-ETSI}-REL & Komponenten replizieren, Zustand handhaben, Ausfälle abfedern & Zuverlässigkeit, Skalierung, Failover \\
\hline
\end{tabular}
\\[1em]
\footnotesize{\textit{Quellen:} \textcite{Antonova2025} (Asterisk-\gls{gls-PABX}, Kamailio, \gls{gls-RTP}-Proxy, Datenbank, ConfigMaps/Geheimnisse, Load Balancer, Monitoring/Logging), \textcite{ETSI-REL} (Zuverlässigkeit/Skalierung als Architekturprinzipien), \textcite{Todorov2025LLMDispatcher} \gls{gls-PABX}, \gls{gls-STT}/\gls{gls-TTS}, \gls{gls-LLM}-Kern, externe \gls{gls-API}, Fallback))}
\end{table}
\FloatBarrier

\noindent \textit{Anmerkung:} Tabelle~\ref{tab:vnf_components} fasst die Komponenten und Architekturprinzipien aus Antonova, \gls{gls-ETSI}-REL und Todorov zusammen. Antonova et al. beschreiben eine Kubernetes-basierte Asterisk-Architektur mit Kamailio, PostgreSQL/MySQL, Load Balancer und Monitoring-Tools. Todorov ergänzt eine produktive Dispatcher-Architektur mit \gls{gls-STT}/\gls{gls-TTS}, \gls{gls-LLM}-Kern, externen \gls{gls-API}-Schnittstellen und Fallback auf menschliche Operatoren. \gls{gls-ETSI}-REL liefert die Prinzipien für Zuverlässigkeit, Skalierung und Failover.

\subsubsection{Abgrenzung: Zielarchitektur und Prototyp-Implementierung}

Die beschriebene Zielarchitektur stellt ein Produktiv-Referenzmodell dar, das auf etablierten Standards und praktischen Validierungen basiert. Die zentralen Komponenten und Prinzipien sind in Tabelle~\ref{tab:vnf_components} zusammengefasst. Der Prototyp dieser Arbeit implementiert ein reduziertes Single-Host-Subset dieser Zielarchitektur.

Kapitel~\ref{chap:implementierung} beschreibt die konkrete Implementierung dieses reduzierten Prototyps und deren Abweichungen von der Zielarchitektur.

\section{Sicherheitskonzept}

Der Prototyp orientiert sich an \gls{gls-ETSI}-Sicherheitsstandards \parencite{ETSI-REL}, implementiert jedoch ein reduziertes Sicherheitsmodell für Entwicklungs- und Testzwecke. Daher unterscheidet sich die Darstellung speziell zwischen Prototyp-Anforderungen und produktiven Sicherheitsma\ss nahmen nach \gls{gls-RFC}- und \gls{gls-ETSI}-Standards.

\subsection{Sicherheitsarchitektur}

\textbf{Prototyp-Modell:}\\
Die folgenden Sicherheitsprinzipien orientieren sich an etablierten Standards (wie \gls{gls-ETSI}-REL und relevanten \gls{gls-RFC}, wurden jedoch spezifisch für den Prototyp angepasst und reduziert:

\begin{itemize}
  \item \textbf{Netzwerk-Isolation:} containerbasierte Segmentierung zur Trennung von Komponenten und Minimierung der Angriffsfläche
  \item \textbf{Zugriffskontrolle:} tokenbasierte Authentifizierung mit definierten Berechtigungen und zeitlich begrenzter Gültigkeit
  \item \textbf{Service-Authentifizierung:} Authentifizierung zwischen internen Services über kryptographische Verfahren
  \item \textbf{Input-Validierung:} Prüfung und Validierung aller Eingabeparameter zur Vermeidung von Injection-Angriffen
  \item \textbf{Nachvollziehbarkeit:} Logging-Mechanismen zur Protokollierung sicherheitsrelevanter Ereignisse
  \item \textbf{Konfigurationsmanagement:} Externe Verwaltung sensibler Daten ohne Einbettung in Quellcode
  \item \textbf{Transportverschlüsselung:} Verzicht auf \gls{gls-TLS}/\gls{gls-SRTP} im Prototyp zugunsten reduzierter Komplexität
\end{itemize}
\textbf{Produktiv-Anforderungen nach Standards:}\\
Produktive Systeme müssen nach \gls{gls-ETSI}-REL \parencite{ETSI-REL} und RFC-Standards folgende Ma\ss nahmen implementieren:

\begin{itemize}
  \item \textbf{\gls{gls-SIP}/\gls{gls-TLS}:} Port 5061 mit \gls{gls-TLS} 1.3 (RFC 8446 \parencite{rfc8446}, RFC 5630 \parencite{rfc5630})
  \item \textbf{\gls{gls-SRTP}:} Verschlüsselte \gls{gls-RTP}-Streams (RFC 3711 \parencite{rfc3711})
  \item \textbf{\gls{gls-DTLS}-\gls{gls-SRTP}:} Für \gls{gls-WebRTC}-Clients (RFC 5764 \parencite{rfc5764})
  \item \textbf{\gls{gls-mTLS}:} zertifikatbasierte Authentifizierung zwischen Services \parencite{rfc8446}
  \item \textbf{\gls{gls-RBAC}:} Rollenbasierte Zugriffskontrolle für administrative Funktionen \parencite{KubernetesDocs}
\end{itemize}

\subsection{Netzwerk-Segmentierung}

Die Architektur folgt dem Prinzip der funktionalen Trennung nach \gls{gls-ETSI}-REL \parencite{ETSI-REL}:

\begin{table}[H]
\centering
\caption[Funktionale Netzwerk-Segmentierung]{Funktionale Netzwerk-Segmentierung nach \gls{gls-ETSI}-Standards}
\label{tab:network_segmentation}
\small
\begin{tabular}{|>{\raggedright\arraybackslash}p{2.5cm}|>{\raggedright\arraybackslash}p{4.5cm}|>{\raggedright\arraybackslash}p{5cm}|}
\hline
\textbf{Zone} & \textbf{Komponenten} & \textbf{Schutzma\ss nahmen} \\
\hline
Signalisierung & \gls{gls-SIP}-Gateway, \gls{gls-NAT}-Traversierung & \gls{gls-TLS} 1.3, \gls{gls-SIP} (RFC 5630) \\
\hline
Medien & \gls{gls-RTP}-Relay, \gls{gls-WebRTC} & \gls{gls-SRTP}, \gls{gls-DTLS} (RFC 3711, 5764) \\
\hline
Anwendung & \gls{gls-KI}-Dienste, \gls{gls-LLM}, \gls{gls-ASR}/\gls{gls-TTS} & \gls{gls-TLS} 1.3 (RFC 8446), \gls{gls-API}-Verschlüsselung \\
\hline
Datenhaltung & Konfiguration, State-Management & Netzwerk-Isolation, Zugriffskontrolle angelehnt an \gls{gls-ITU-T} \\
\hline
\end{tabular}
\\[1em]
\footnotesize{\textit{Quelle:} \textcite{ETSI-REL}}
\end{table}
\FloatBarrier

Die folgende Zielarchitektur wurde auf der Grundlage etablierter Architekturprinzipien und unter Einbeziehung relevanter Literatur abgeleitet \parencite{ETSI-NFV,Antonova2025,Todorov2025LLMDispatcher}.

\begin{figure}[H]
\centering
\includegraphics[width=0.85\textwidth,height=0.85\textheight,keepaspectratio]{images/architektur.png}
\caption{Architektur-Übersicht: Eigene Darstellung mit PlantUML
(\url{https://plantuml.com/}) in Anlehnung an relevante Literatur.}
\label{fig:architektur}
\end{figure}

Die konkrete Implementierung dieser Architektur-Prinzipien wird in Kapitel~\ref{chap:implementierung} beschrieben.

\subsection{Bereitstellungsmodelle und Deployment}

Bei der Systemarchitektur stehen drei grundlegende Modelle für eine Bereitstellung zur Verfügung: Das \textbf{monolithische Modell} fasst alle Funktionen in einer einzigen Container-Instanz zusammen, was schnelle Prototypentwicklung ermöglicht, aber Skalierungsprobleme mit sich bringt. Das \textbf{vollständig dezentralisierte Modell} zerlegt die Anwendung in viele einzelne, unabhängige Mikroservices, bietet höhere Flexibilität und Ausfallsicherheit, erfordert aber umfangreiche Orchestrierungs- und Verwaltungsma\ss nahmen. Der \textbf{hybride Ansatz} stellt einen Mittelweg dar: Er isoliert nur die kritischsten Funktionen als eigenständige Services, während unkritische Funktionen weiterhin zusammengefasst bleiben. Der vorliegende Prototyp folgt dieser hybriden Strategie, um eine schrittweise Entkopplung von Komponenten ohne umfassende Neuarchitekturierung zu ermöglichen \parencite{Antonova2025}.

Im folgenden Kapitel~\ref{chap:marktanalyse} wird nun genauer auf einige Anbieter im Bereich der \gls{gls-KI}-basierten Telefonie eingegangen. Es werden exemplarisch sieben Anbieter betrachtet, die unterschiedlichste Architekturen nutzen, diverse Zielgruppen adressieren und verschiedene Kernfeatures verwenden. Au\ss erdem werden ausführlich ihre Preismodelle dargestellt und mögliche anfallende Kosten analysiert. Zusätzlich werden wichtige datenschutzrechtliche Aspekte diskutiert und abgeleitete Ma\ss nahmen empfohlen.
